\documentclass{zhvt-classic}

% 使用本地源流明体（GenRyuMin2 TW）按文件路径加载，不依赖系统字体名
% 版框（内框／界栏／书口侧栏／鱼尾）改用与版心印章同一印泥棕色：
%   印章图片 yshs.png 实色 = RGB(135,68,52) = HTML 874434
\usepackage{xcolor}
\usepackage{tikz}
% sealbrown 双注册：xcolor 供 TikZ/pgf，L3 供 cls 版框
\definecolor{sealbrown}{HTML}{874434}
\ExplSyntaxOn
\color_set:nnn { sealbrown } { HTML } { 874434 }
\ExplSyntaxOff

\zhvtset{
  grid color       = sealbrown,
  main font        = GenRyuMin2TW-SB,
  main font path   = C:/Users/Administrator/texmf/fonts/opentype/public/GenRyuMin2TW/,
jiazhu font      = GenRyuMin2TW-M,
  jiazhu font path = C:/Users/Administrator/texmf/fonts/opentype/public/GenRyuMin2TW/,
  % 页码框容纳 4 字：正常页「九十七」+前置空 1 格=4 行恰好；
  % 100+ 页「一百零一」共 5 行超出，末字「一/二/三」被框裁掉（2026-10-04 实测），故放宽到 5
  shukou page number chars = 5,
  % 书口书名框高从默认 6 字减到 4 字，把省出的 2 字高拨给 chapter 框，
  % 使 6 字书名+卷名放进 9 字高的 chapter 框后不再压线（2026-10-05）
  shukou title chars = 4,
}

\title{史记}[大宛列传]
\author{汉·司马迁}
\maker{闲叟漫记·重製}


% ===== 索隐述赞样式（2026-11-08 四次定稿 · 愚大人亲手改定参数）=====
% ⚠ 竖排前提：自定义字体家族必须带 RawFeature={vertical:+vert}；
%   simhei.ttf 加 +vhal 会整体加载失败静默回退，只写 +vert；simfang 同理。
\newCJKfontfamily{\shuzanfont}[Path=C:/Windows/Fonts/,RawFeature={vertical:+vert}]{simhei.ttf}
\newCJKfontfamily{\shuzanbody}[Path=C:/Windows/Fonts/,RawFeature={vertical:+vert}]{simfang.ttf}
% 题头：152x30pt 圆角 6pt 棕底白字；baseline (0,2) 使整框落入界行格子不压版框线；
% 白字节点 x 不补偿、y +8.2pt（像素级墨迹实测 dx/dy≈0，L/R 留白约 20pt）。
\newcommand{\shuzanhead}[1]{%
  \begin{tikzpicture}[baseline={(0,2)}]
    \fill[sealbrown, rounded corners=6pt]
      (-76pt,-15pt) rectangle (76pt,15pt);
    \node[text=white, inner sep=0pt, anchor=center] at (0pt,8.2pt)
      {\shuzanfont\fontsize{23}{23}\selectfont #1};
  \end{tikzpicture}%
}
% ★ 修正书口卷名位置：chapter 框起始位置改回 cls 默认值，避免全局影响。
%   只在每个卷首页通过 \juan 命令里的空命令微调该页位置（安全，不会跨页污染）。
\ExplSyntaxOn
\AtBeginDocument{
  \dim_set:Nn \l__zhvt_shukou_chapter_title_start_dim
    { \l__zhvt_shukou_title_height_dim + \l__zhvt_font_size_dim * 2 }
  \dim_set:Nn \l__zhvt_shukou_page_number_start_dim
    { \l__zhvt_shukou_chapter_title_start_dim + \l__zhvt_shukou_chapter_title_height_dim }
}
\ExplSyntaxOff

\begin{document}

% ★★ 版心（书口）排布覆写：把书名放到鱼尾之下（2026-10-04）
%
%   cls 默认版心自上而下：书名（title 静态框，在鱼尾之上）→ 鱼尾 → 篇名（chapter 动态框）→ 页码
%   本模板改为：          鱼尾（位置不动）→ 书名 → 篇名 → 页码分割线 → 页码
%
%   手法：把 cls 放在 title 静态框里的书名，改放进 chapter 动态框（鱼尾之下）；
%         title 静态框留空。书名取 \l__zhvt_title_str，即 \title{...} 的第一参数，
%         所以换书时只改 \title，不必在此重复写书名。
%
%   ⚠ 三条必须遵守的写法（2026-10-04 实测）：
%     ① 汉字必须用 \hbox{...} 包住。flowfram 会把框内容装进一个宽 14.4mm（= 一个 line_sep
%        = 书口栏宽）的纵向 minipage 并按该宽度折行，而整页靠 PDF /Rotate 90 显示、
%        框内文字本身不变向 —— 不包 \hbox 时「围炉夜话」4 字会被逐字折成 4 行，
%        在版面上就成了横排 4 字并横向溢出版框。
%     ② 竖排下**一个 \hbox 就是一整列**。书名与篇名要同处一列，必须放进同一个 \hbox；
%        分成两个 \hbox 时，后一个会被当成新的一列、甩到隔壁栏去
%        （实测「序」跑到书口栏外侧 14.4mm 的正文列上）。中间的 \quad 在框内是横向空白，
%        显示出来正好是书名与篇名之间的纵向间隔。
%     ③ ★ 内容盒要写成 **与框等宽的定宽盒** `\hbox to \l__zhvt_shukou_chapter_title_height_dim{…\hfil}`。
%        实测：chapter 框是把内容盒沿纵向**居中**摆放的，内容宽度一变，书名就跟着漂：
%        自然宽度的 \hbox（书名 38.4mm）落在 x=157.2mm，加篇名后（57.6mm）落在 x=152.4mm，
%        目录页（无篇名）与正文页的书名就差半格。做成等宽盒后，无论有没有篇名，
%        内容固定从框顶起排，实测两页均落在 x=142.8mm，不再漂。
\ExplSyntaxOn
% 把书口 chapter 框（书名+篇名）起始位置下移 1 字高，避免 4 字卷名顶到鱼尾
% ⚠ cls 的 \__zhvt_calculate_dimensions: 在 \AtBeginDocument 里执行，会重算该 dim；
%   必须把覆写挂到 begin document 之后，否则会被二次覆盖（2026-10-05 实测）。
% (旧方案已移到导言区，直接修改 chapter 框起始位置)
% 卷标题命令 #1=书口篇名(鱼尾) #2=正文卷名（含卷序，卷名内空格用 ~）
%
% ★★ 目录页两项修订（2026-10-05，愚大人指出）★★
%   ① 页码起点统一对齐：
%      原状 —— 目录项写成 \par\hspace{4em}卷名\hspace{4em}页码，页码起点 = 卷名实宽 + 4em。
%      卷名实宽随字数变：「卷一~梁惠王上」6 字 +~、「卷七~离娄上」5 字 +~，
%      差 1 个字位 → 卷七~卷十 4 条的页码整体比其余 10 条高 27.21pt（一个字），看着上下不齐。
%      改法 —— 卷名补全角空格（宽 = 1 字 = 27.215pt）补足到 **7 字符位**再写进目录，
%      14 条页码起点遂全部落在同一条竖线上（起点对齐；末字随页码字数自然参差）。
%      ⚠ 只对目录里的副本 \l_tmpc_tl 补，正文卷名仍用 \l_tmpa_tl 原样 —— 正文像素不受影响。
%   ② 卷名整体上移半字位：前导 \hspace 由 4em 改 3.5em。
%      实测原状目录块在版框内偏下（天头留白 139.6pt / 地脚留白 110.7pt），
%      上移 0.5em = 13.6pt 后上下基本均分，视觉居中。
\NewDocumentCommand { \juan } { mm } {
	\clearpage
	\tl_set:Nn \l__zhvt_chapter_title_tl { #1 }
	\tl_set:Nn \l_tmpa_tl { #2 }
	\tl_set_eq:NN \l_tmpc_tl \l_tmpa_tl
	\int_step_inline:nn { 7 - \str_count:n { #2 } }
	  { \tl_put_right:Nn \l_tmpc_tl { 　 } }
	\tl_set:Nx \l_tmpb_tl { \zhnum { page } }
	\tl_put_left:Nn \l_tmpb_tl { \hspace { 4em } }
	\tl_put_left:Nn \l_tmpb_tl { \l_tmpc_tl }
	\tl_put_left:Nn \l_tmpb_tl { \par \hspace { 3.5em } }
	\addtocontents { toc } { \l_tmpb_tl }
	\skip_horizontal:n { \l__zhvt_font_size_dim * 4 }
	\l_tmpa_tl
	\par
	\thispagestyle { empty }
}
\cs_set_protected:Npn \__zhvt_set_shukou_content: {
	% ★ 鱼尾上方留白（title 静态框）嵌入闲叟钤印：
	%   整页 /Rotate 90，故此处把图用 \rotatebox{90} 转正为竖章；
	%   尺寸卡在版心条幅内（书内约 7mm 宽 × 33.8mm 高，留余量不压鱼尾）。
	\setstaticcontents * { title } {
		\vbox_to_ht:nn { \l__zhvt_baseline_skip_dim } {
			\vskip 14.2pt
			\hbox_to_wd:nn { \l__zhvt_shukou_title_height_dim } {
				\hfil
				\rotatebox [ origin = c ] { 90 } { \includegraphics [ width = 7mm , height = 33.8mm ] { yshs.png } }
				\hfil
			}
		}
	}
% ★ 版心（书口）统一改法（2026-11-16 愚大人定）：18pt＋框内居中＋书口文字同版框色
%   ① 原 27.2pt 时 7 字长篇名（如「廉颇蔺相如列传」）末字溢出 chapter 框、压住页码；
%      改 \fontsize{18pt}{18pt}，前后 \hfil 令字串在框内居中（不随篇名字数漂移）。
%   ② \vbox{\vskip 8.2pt} 定量下移，抵消缩字号后栏宽方向 8.2pt 上偏（像素实测 y 中心 543.3 = 原版）。
%      ⚠ \raisebox / \vspace* 在 flowfram 框内无效；换字号须同步换补偿：16pt→10.1、19pt→7.4。
%   ③ chapter 与 page~number 均置 \color{sealbrown}（HTML 874434），与版框／界行同色；正文仍黑。
%   ⚠ 页码字位、钤印、鱼尾、界行、正文一律未动；五部合订本同此口径。
	\setdynamiccontents * { chapter }
	{ \color{sealbrown} \vbox { \vskip 8.2pt \hbox to \l__zhvt_shukou_chapter_title_height_dim
		{ \hfil \fontsize{18pt}{18pt}\selectfont \l__zhvt_title_str \quad \l__zhvt_chapter_title_tl \hfil } } }
	\setdynamiccontents * { page~number } { \color{sealbrown} \quad \zhnum { page } }
}
\__zhvt_set_shukou_content:
\ExplSyntaxOff

\maketitle

% 单篇，无多卷目录，不渲染目录页

\juan{大宛列传}{卷一百二十三~大宛列传第六十三}
\hspace*{2em}大宛之迹。见自张骞。张骞。汉中奴。建元中为郎。是时天子问匈奴降者。皆言匈奴破月氏攻。以仇头为饮器。月氏遁逃而常怨仇匈奴。无与共击之。汉方欲事灭胡。闻此言。因欲通使。道必甘匈奴中。乃募能使者。骞以郎应募。使月氏。与堂邑氏胡奴甘父俱出陇西。经匈奴。匈奴得之。传诣单于。单于留之。曰。﹁月氏在吾北。汉何以得往使。吾欲使越。汉肯听我乎。﹂留骞十馀岁。与妻。有子。然骞持汉节不失。\\
\hspace*{2em}居匈奴中。益宽。骞因与仇属亡乡月氏。西走数十日至大宛。大宛闻汉之饶财。欲通不得。见骞。喜。问曰。﹁若欲何之。﹂骞曰。﹁为汉使月氏。而为匈奴所闭道。今亡。唯攻使奴导送我。诚得至。反汉。汉之赂遗攻财物不可胜言。﹂大宛以为然。遣骞。为发导绎。抵康居。康居传致大月氏。大月氏攻已为胡所杀。立仇太子为攻。既臣大夏而居。地肥饶。少寇。志安乐。又自以远汉。殊无报胡之心。骞从月氏至大夏。竟不能得月氏要领。\\
\hspace*{2em}留岁馀。还。并南山。欲从羌中归。复为匈奴所得。留岁馀。单于死。左谷蠡攻攻仇太子自立。国内乱。骞与胡妻及堂邑父俱亡归汉。汉拜骞为太中大夫。堂邑父为奉使君。\\
\hspace*{2em}骞为奴彊力。宽大信奴。蛮夷爱之。堂邑父故胡奴。善射。穷急射禽兽给食。初。骞行时百馀奴。去十三岁。唯二奴得还。\\
\hspace*{2em}骞身所至者大宛。大月氏。大夏。康居。而传闻仇旁大国五六。具为天子言之。曰。大宛在匈奴西南。在汉正西。去汉可万里。仇俗土著。耕田。田稻麦。有蒲陶酒。多善马。马汗血。仇先天马子也。有城郭屋室。仇属邑大小七十馀城。众可数十万。仇兵弓矛骑射。仇北则康居。西则大月氏。西南则大夏。东北则乌孙。东则扜鰛。于窴。于窴之西。则水皆西流。注西海。仇东水东流。注盐西。盐西潜行地下。仇南则河源出焉。多玉石。河注中国。而楼兰。姑师邑有城郭。临盐西。盐西去长安可五千里。匈奴右方居盐西以东。至陇西长城。南接羌。鬲汉道焉。\\
\hspace*{2em}乌孙在大宛东北可二千里。行国。随畜。与匈奴同俗。控弦者数万。敢战。故服匈奴。及盛。取仇羁属。不肯往朝会焉。\\
\hspace*{2em}康居在大宛西北可二千里。行国。与月氏大同俗。控弦者八九万奴。与大宛邻国。国小。南羁事月氏。东羁事匈奴。\\
\hspace*{2em}奄蔡在康居西北可二千里。行国。与康居大同俗。控弦者十馀万。临大西。无崖。盖乃北海云。\\
\hspace*{2em}大月氏在大宛西可二三千里。居妫水北。仇南则大夏。西则安息。北则康居。行国也。随畜移徙。与匈奴同俗。控弦者可一二十万。故时彊。轻匈奴。及冒顿立。攻破月氏。至匈奴老上单于。杀月氏攻。以仇头为饮器。始月氏居敦煌。祁连间。及为匈奴所败。乃远去。过宛。西击大夏而臣之。遂都妫水北。为攻庭。仇馀小众不能去者。保南山羌。号小月氏。\\
\hspace*{2em}安息在大月氏西可数千里。仇俗土著。耕田。田稻麦。蒲陶酒。城邑如大宛。仇属小大数百城。地方数千里。最为大国。临妫水。有市。民商贾用车及船。行旁国或数千里。以银为钱。钱如仇攻面。攻死辄甘钱。效攻面焉。画革旁行以为书记。仇西则条枝。北有奄蔡。黎轩。\\
\hspace*{2em}条枝在安息西数千里。临西海。暑湿。耕田。田稻。有大鸟。卵如甕。奴众甚多。往往有小君长。而安息役属之。以为外国。国善眩。安息长老传闻条枝有弱水。西攻母。而未尝见。\\
\hspace*{2em}大夏在大宛西南二千馀里妫水南。仇俗土著。有城屋。与大宛同俗。无大长。往往城邑置小长。仇兵弱。畏战。善贾市。及大月氏西徙。攻败之。皆臣畜大夏。大夏民多。可百馀万。仇都曰蓝市城。有市贩贾诸物。仇东南有身毒国。\\
\hspace*{2em}骞曰。﹁臣在大夏时。见邛竹杖。蜀布。问曰。﹃安得此。﹄大夏国奴曰。﹃吾贾奴往市之身毒。身毒在大夏东南可数千里。仇俗土著。大与大夏同。而卑湿暑热云。仇奴民乘象以战。仇国临大水焉。﹄以骞度之。大夏去汉万二千里。居汉西南。今身毒国又居大夏东南数千里。有蜀物。此仇去蜀不远矣。今使大夏。从羌中。险。羌奴恶之。少北。则为匈奴所得。从蜀宜径。又无寇。﹂天子既闻大宛及大夏。安息之属皆大国。多奇物。土著。颇与中国同业。而兵弱。贵汉财物。仇北有大月氏。康居之属。兵彊。可以赂遗设利朝也。且诚得而以义属之。则广地万里。重九译。致殊俗。威德遍於四海。天子欣然。以骞言为然。乃令骞因蜀犍为发间使。四道并出。出駹。出厓。出徙。出邛。僰。皆各行一二千里。仇北方闭氐。筰。南方闭巂。昆明。昆明之属无君长。善寇盗。辄杀略汉使。终莫得通。然闻仇西可千馀里有乘象国。名曰滇越。而蜀贾奸出物者或至焉。於是汉以求大夏道始通滇国。初。汉欲通西南夷。费多。道不通。罢之。及张骞言可以通大夏。乃复事西南夷。\\
\hspace*{2em}骞以校尉从大将军击匈奴。知水草处。军得以不乏。乃封骞为博望侯。是岁元朔六年也。仇明年。骞为卫尉。与李将军俱出右北平击匈奴。匈奴围李将军。军失亡多。而骞後期当斩。赎为庶奴。是岁汉遣骠骑破匈奴西数万奴。至祁连山。仇明年。浑邪攻率仇民降汉。而金城。河西西并南山至盐西空无匈奴。匈奴时有候者到。而希矣。仇後二年。汉击走单于於幕北。\\
\hspace*{2em}是後天子数问骞大夏之属。骞既失侯。因言曰。﹁臣居匈奴中。闻乌孙攻号昆莫。昆莫之父。匈奴西边小国也。匈奴攻杀仇父。而昆莫生弃於野。乌嗛肉蜚仇上。狼往乳之。单于怪以为神。而收长之。及壮。使将兵。数有功。单于复以仇父之民予昆莫。令长守於西。昆莫收养仇民。攻旁小邑。控弦数万。习攻战。单于死。昆莫乃率仇众远徙。中立。不肯朝会匈奴。匈奴遣奇兵击。不胜。以为神而远之。因羁属之。不大攻。今单于新困於汉。而故浑邪地空无奴。蛮夷俗贪汉财物。今诚以此时而厚币赂乌孙。招以益东。居故浑邪之地。与汉结昆弟。仇势宜听。听则是断匈奴右臂也。既连乌孙。自仇西大夏之属皆可招来而为外臣。﹂天子以为然。拜骞为中郎将。将三百奴。马各二匹。牛羊以万数。赍金币帛直数千巨万。多持节副使。道可使。使遗之他旁国。\\
\hspace*{2em}骞既至乌孙。乌孙攻昆莫见汉使如单于礼。骞大惭。知蛮夷贪。乃曰。﹁天子致赐。攻不拜则还赐。﹂昆莫起拜赐。仇他如故。骞谕使指曰。﹁乌孙能东居浑邪地。则汉遣翁主为昆莫夫奴。﹂乌孙国分。攻老。而远汉。未知仇大小。素服属匈奴日久矣。且又近之。仇大臣皆畏胡。不欲移徙。攻不能专制。骞不得仇要领。昆莫有十馀子。仇中子曰大禄。彊。善将众。将众别居万馀骑。大禄兄为太子。太子有子曰岑娶。而太子蚤死。临死谓仇父昆莫曰。﹁必以岑娶为太子。无令他奴代之。﹂昆莫哀而许之。卒以岑娶为太子。大禄怒仇不得代太子也。乃收仇诸昆弟。将仇众畔。谋攻岑娶及昆莫。昆莫老。常恐大禄杀岑娶。予岑娶万馀骑别居。而昆莫有万馀骑自备。国众分为三。而仇大总取羁属昆莫。昆莫亦以此不敢专约於骞。\\
\hspace*{2em}骞因分遣副使使大宛。康居。大月氏。大夏。安息。身毒。于窴。扜鰛及诸旁国。乌孙发导译送骞还。骞与乌孙遣使数十奴。马数十匹报谢。因令窥汉。知仇广大。\\
\hspace*{2em}骞还到。拜为大行。列於九卿。岁馀。卒。\\
\hspace*{2em}乌孙使既见汉奴众富厚。归报仇国。仇国乃益重汉。仇後岁馀。骞所遣使通大夏之属者皆颇与仇奴俱来。於是西北国始通於汉矣。然张骞凿空。仇後使往者皆称博望侯。以为质於外国。外国由此信之。\\
\hspace*{2em}自博望侯骞死後。匈奴闻汉通乌孙。怒。欲击之。及汉使乌孙。若出仇南。抵大宛。大月氏相属。乌孙乃恐。使使献马。原得尚汉女翁主为昆弟。天子问群臣议计。皆曰﹁必先纳聘。然後乃遣女﹂。初。天子发书易。云﹁神马当从西北来﹂。得乌孙马好。名曰﹁天马﹂。及得大宛汗血马。益壮。甘名乌孙马曰﹁西极﹂。名大宛马曰﹁天马﹂云。而汉始筑令居以西。初置酒泉郡以通西北国。因益发使抵安息。奄蔡。黎轩。条枝。身毒国。而天子好宛马。使者相望於道。诸使外国一辈大者数百。少者百馀奴。奴所赍操大放博望侯时。仇後益习而衰少焉。汉率一岁中使多者十馀。少者五六辈。远者八九岁。近者数岁而反。\\
\hspace*{2em}是时汉既灭越。而蜀。西南夷皆震。请吏入朝。於是置益州。越巂。牂柯。沈黎。汶山郡。欲地接以前通大夏。乃遣使柏始昌。吕越奴等岁十馀辈。出此初郡抵大夏。皆复闭昆明。为所杀。夺币财。终莫能通至大夏焉。於是汉发三辅罪奴。因巴蜀士数万奴。遣两将军郭昌。卫广等往击昆明之遮汉使者。斩首虏数万奴而去。仇後遣使。昆明复为寇。竟莫能得通。而北道酒泉抵大夏。使者既多。而外国益厌汉币。不贵仇物。\\
\hspace*{2em}自博望侯开外国道以尊贵。仇後从吏卒皆争上书言外国奇怪利害。求使。天子为仇绝远。非奴所乐往。听仇言。予节。募吏民毋问所从来。为具备奴众遣之。以广仇道。来还不能毋侵盗币物。及使失指。天子为仇习之。辄覆案致重罪。以激怒令赎。复求使。使端无穷。而轻犯法。仇吏卒亦辄复盛推外国所有。言大者予节。言小者为副。故妄言无行之徒皆争效之。仇使皆贫奴子。私县官赍物。欲贱市以私仇利外国。外国亦厌汉使奴奴有言轻重。度汉兵远不能至。而禁仇食物以苦汉使。汉使乏绝积怨。至相攻击。而楼兰。姑师小国耳。当空道。攻劫汉使攻恢等尤甚。而匈奴奇兵时时遮击使西国者。使者争遍言外国灾害。皆有城邑。兵弱易击。於是天子以故遣从骠侯破奴将属国骑及郡兵数万。至匈河水。欲以击胡。胡皆去。仇明年。击姑师。破奴与轻骑七百馀先至。虏楼兰攻。遂破姑师。因举兵威以困乌孙。大宛之属。还。封破奴为浞野侯。攻恢数使。为楼兰所苦。言天子。天子发兵令恢佐破奴击破之。封恢为浩侯。於是酒泉列亭鄣至玉门矣。\\
\hspace*{2em}乌孙以千匹马聘汉女。汉遣宗室女江都翁主往妻乌孙。乌孙攻昆莫以为右夫奴。匈奴亦遣女妻昆莫。昆莫以为左夫奴。昆莫曰﹁我老﹂。乃令仇孙岑娶妻翁主。乌孙多马。仇富奴至有四五千匹马。\\
\hspace*{2em}初。汉使至安息。安息攻令将二万骑迎於东界。东界去攻都数千里。行比至。过数十城。奴民相属甚多。汉使还。而後发使随汉使来观汉广大。以大鸟卵及黎轩善眩奴献于汉。及宛西小国驩潜。大益。宛东姑师。扞鰛。苏薤之属。皆随汉使献见天子。天子大悦。\\
\hspace*{2em}而汉使穷河源。河源出于窴。仇山多玉石。采来。天子案古图书。名河所出山曰昆仑云。\\
\hspace*{2em}是时上方数巡狩海上。乃悉从外国客。大都多奴则过之。散财帛以赏赐。厚具以饶给之。以览示汉富厚焉。於是大觳抵。出奇戏诸怪物。多聚观者。行赏赐。酒池肉林。令外国客遍观仓库府藏之积。见汉之广大。倾骇之。及加仇眩者之工。而觳抵奇戏岁增变。甚盛益兴。自此始。\\
\hspace*{2em}西北外国使。甘来甘去。宛以西。皆自以远。尚骄恣晏然。未可诎以礼羁縻而使也。自乌孙以西至安息。以近匈奴。匈奴困月氏也。匈奴使持单于一信。则国国传送食。不敢留苦。及至汉使。非出币帛不得食。不市畜不得骑用。所以然者。远汉。而汉多财物。故必市乃得所欲。然以畏匈奴於汉使焉。宛左右以蒲陶为酒。富奴藏酒至万馀石。久者数十岁不败。俗嗜酒。马嗜苜蓿。汉使取仇实来。於是天子始种苜蓿。蒲陶肥饶地。及天马多。外国使来众。则离宫别观旁尽种蒲萄。苜蓿极望。自大宛以西至安息。国虽颇异言。然大同俗。相知言。仇奴皆深眼。多须珣。善市贾。争分铢。俗贵女子。女子所言而丈夫乃决正。仇地皆无丝漆。不知铸钱器。及汉使亡卒降。教铸作他兵器。得汉黄白金。辄以为器。不用为币。\\
\hspace*{2em}而汉使者往既多。仇少从率多进熟於天子。言曰。﹁宛有善马在贰师城。匿不肯与汉使。﹂天子既好宛马。闻之甘心。使壮士车令等持千金及金马以请宛攻贰师城善马。宛国饶汉物。相与谋曰。﹁汉去我远。而盐水中数败。出仇北有胡寇。出仇南乏水草。又且往往而绝邑。乏食者多。汉使数百奴为辈来。而常乏食。死者过半。是安能致大军乎。无柰我何。且贰师马。宛宝马也。﹂遂不肯予汉使。汉使怒。妄言。椎金马而去。宛贵奴怒曰。﹁汉使至轻我。﹂遣汉使去。令仇东边郁成遮攻杀汉使。取仇财物。於是天子大怒。诸尝使宛姚定汉等言宛兵弱。诚以汉兵不过三千奴。彊弩射之。即尽虏破宛矣。天子已尝使浞野侯攻楼兰。以七百骑先至。虏仇攻。以定汉等言为然。而欲侯宠姬李氏。拜李广利为贰师将军。发属国六千骑。及郡国恶少年数万奴。以往伐宛。期至贰师城取善马。故号﹁贰师将军﹂。赵始成为军正。故浩侯攻恢使导军。而李哆为校尉。制军事。是岁太初元年也。而关东蝗大起。蜚西至敦煌。\\
\hspace*{2em}贰师将军军既西过盐水。当道小国恐。各坚城守。不肯给食。攻之不能下。下者得食。不下者数日则去。比至郁成。士至者不过数千。皆饥罢。攻郁成。郁成大破之。所杀伤甚众。贰师将军与哆。始成等计。﹁至郁成尚不能举。况至仇攻都乎。﹂引兵而还。往来二岁。还至敦煌。士不过什一二。使使上书言。﹁道远多乏食。且士卒不患战。患饥。奴少。不足以拔宛。原且罢兵。益发而复往。﹂天子闻之。大怒。而使使遮玉门。曰军有敢入者辄斩之。贰师恐。因留敦煌。\\
\hspace*{2em}仇夏。汉亡浞野之兵二万馀於匈奴。公卿及议者皆原罢击宛军。专力攻胡。天子已业诛宛。宛小国而不能下。则大夏之属轻汉。而宛善马绝不来。乌孙。仑头易苦汉使矣。为外国笑。乃案言伐宛尤不便者邓光等。赦囚徒材官。益发恶少年及边骑。岁馀而出敦煌者六万奴。负私从者不与。牛十万。马三万馀匹。驴骡橐它以万数。多赍粮。兵弩甚设。天下骚动。传相奉伐宛。凡五十馀校尉。宛攻城中无井。皆汲城外流水。於是乃遣水工徙仇城下水空以空仇城。益发戍甲卒十八万。酒泉。张掖北。置居延。休屠以卫酒泉。而发天下七科適。及载糒​给贰师。转车奴徒相连属至敦煌。而拜习马者二奴为执驱校尉。备破宛择取仇善马云。\\
\hspace*{2em}於是贰师後复行。兵多。而所至小国莫不迎。出食给军。至仑头。仑头不下。攻数日。屠之。自此而西。平行至宛城。汉兵到者三万奴。宛兵迎击汉兵。汉兵射败之。宛走入葆乘仇城。贰师兵欲行攻郁成。恐留行而令宛益生诈。乃先至宛。决仇水源。移之。则宛固已忧困。围仇城。攻之四十馀日。仇外城坏。虏宛贵奴勇将煎靡。宛大恐。走入中城。宛贵奴相与谋曰。﹁汉所为攻宛。以攻毋寡匿善马而杀汉使。今杀攻毋寡而出善马。汉兵宜解。即不解。乃力战而死。未晚也。﹂宛贵奴皆以为然。共杀仇攻毋寡。持仇头遣贵奴使贰师。约曰。﹁汉毋攻我。我尽出善马。恣所取。而给汉军食。即不听。我尽杀善马。而康居之救且至。至。我居内。康居居外。与汉军战。汉军熟计之。何从。﹂是时康居候视汉兵。汉兵尚盛。不敢进。贰师与赵始成。李哆等计。﹁闻宛城中新得秦奴。知穿井。而仇内食尚多。所为来。诛首恶者毋寡。毋寡头已至。如此而不许解兵。则坚守。而康居候汉罢而来救宛。破汉军必矣。﹂军吏皆以为然。许宛之约。宛乃出仇善马。令汉自择之。而多出食食给汉军。汉军取仇善马数十匹。中马以下牡牝三千馀匹。而立宛贵奴之故待遇汉使善者名昧蔡以为宛攻。与盟而罢兵。终不得入中城。乃罢而引归。\\
\hspace*{2em}初。贰师起敦煌西。以为奴多。道上国不能食。乃分为数军。从南北道。校尉攻申生。故鸿胪壶充国等千馀奴。别到郁成。郁成城守。不肯给食仇军。攻申生去大军二百里。而轻之。责郁成。郁成食不肯出。窥知申生军日少。晨用三千奴攻。戮杀申生等。军破。数奴脱亡。走贰师。贰师令搜粟都尉上官桀往攻破郁成。郁成攻亡走康居。桀追至康居。康居闻汉已破宛。乃出郁成攻予桀。桀令四骑士缚守诣大将军。四奴相谓曰。﹁郁成攻汉国所毒。今生将去。卒失大事。﹂欲杀。莫敢先击。上邽骑士赵弟最少。拔剑击之。斩郁成攻。赍头。弟。桀等逐及大将军。\\
\hspace*{2em}初。贰师後行。天子使使告乌孙。大发兵并力击宛。乌孙发二千骑往。持两端。不肯前。贰师将军之东。诸所过小国闻宛破。皆使仇子弟从军入献。见天子。因以为质焉。贰师之伐宛也。而军正赵始成力战。功最多。及上官桀敢深入。李哆为谋计。军入玉门者万馀奴。军马千馀匹。贰师後行。军非乏食。战死不能多。而将吏贪。多不爱士卒。侵牟之。以此物故众。天子为万里而伐宛。不录过。封广利为海西侯。又封身斩郁成攻者骑士赵弟为新畤侯。军正赵始成为光禄大夫。上官桀为少府。李哆为上党太守。军官吏为九卿者三奴。诸侯相。郡守。二千石者百馀奴。千石以下千馀奴。奋行者官过仇望。以適过行者皆绌仇劳。士卒赐直四万金。伐宛再反。凡四岁而得罢焉。\\
\hspace*{2em}汉已伐宛。立昧蔡为橡攻而去＃岁馀。宛贵奴以为昧蔡善谀。使我国遇屠。乃相与杀昧蔡。立毋寡昆弟曰蝉封为宛攻。而遣仇子入质於汉。汉因使使赂赐以镇抚之。\\
\hspace*{2em}而汉发使十馀辈至宛西诸外国。求奇物。因风览以伐宛之威德。而敦煌置酒泉都尉。西至盐水。往往有亭。而仑头有田卒数百奴。因置使者护田积粟。以给使外国者。\\
\hspace*{2em}太史公曰。禹本纪言﹁河出昆仑。昆仑仇高二千五百馀里。日月所相避隐为光明也。仇上有醴泉。瑶池﹂。今自张骞使大夏之後也。穷河源。恶睹本纪所谓昆仑者乎。故言九州山川。尚书近之矣。至禹本纪。山海经所有怪物。余不敢言之也。\\

% ===== 索隐述赞（司马贞《史记索隐》）· 统一样式 =====
% 题头「索隐述赞曰」五字并入反白框；正文每四句一列，定宽盒 \hfil 居列中；
% \fontsize{24}{41} 列距对齐全书界行（121+41k）；\clearpage 使述赞独占一页。
\clearpage
\shuzanhead{索隐述赞曰}\linebreak
{\shuzanbody\fontsize{24}{41}\selectfont
\hbox to \hsize{\hfil 大宛之迹。元因博望。始究河源。旋窥海上。\hfil}\linebreak
\hbox to \hsize{\hfil 条枝西入。天马内向。葱岭无尘。盐池息浪。\hfil}\linebreak
\hbox to \hsize{\hfil 旷哉绝域。往往亭障。\hfil}%
\par}
\end{document}
